% ===== PRÉAMBULE CANONIQUE — à coller VERBATIM en tête de CHAQUE .tex =====
\documentclass[11pt,a4paper]{article}
\usepackage[utf8]{inputenc}\usepackage[T1]{fontenc}\usepackage{lmodern}
\usepackage[french]{babel}
\usepackage{amsmath,amssymb,mathtools}\usepackage{icomma}
\usepackage[version=4]{mhchem}          % \ce{...} : équations & formules
\usepackage{siunitx}\sisetup{locale=FR} % \qty{}{}, \si{}, \ang{}
\usepackage{chemfig}                    % \chemfig{...} : structures développées
\usepackage{xcolor}\usepackage{colortbl}
\usepackage{tikz}\usepackage{pgfplots}\pgfplotsset{compat=1.18}
\usepackage[most]{tcolorbox}\usepackage{enumitem}\usepackage{array,booktabs}
\usepackage[a4paper,margin=2.2cm]{geometry}
\usetikzlibrary{arrows.meta,calc,angles,quotes,positioning,patterns,backgrounds,shapes.geometric}
\definecolor{primary}{RGB}{222,49,99}\definecolor{primarydark}{RGB}{150,22,55}
\definecolor{defcol}{RGB}{0,114,178}\definecolor{methcol}{RGB}{52,92,148}
\definecolor{excol}{RGB}{0,135,90}\definecolor{attncol}{RGB}{200,40,55}
\tcbset{boxbase/.style={enhanced,breakable,boxrule=0.9pt,arc=2.5pt,
  left=3mm,right=3mm,top=2.3mm,bottom=2.3mm,fonttitle=\bfseries,coltitle=white}}
% --- boîtes TUTORIEL (noms EXACTS, ne pas renommer) ---
\newtcolorbox{cle}[1][]{boxbase,colback=defcol!6,colframe=defcol,
  title={\textsc{À retenir}\ifx\relax#1\relax\relax\else\textmd{ : \itshape #1}\fi}}
\newtcolorbox{methode}[1][]{boxbase,colback=methcol!7,colframe=methcol,sharp corners,
  title={\textsc{Méthode}\ifx\relax#1\relax\relax\else\textmd{ --- \itshape #1}\fi}}
\newtcolorbox{exemplebox}[1][]{boxbase,colback=excol!5,colframe=excol,
  title={\textsc{Exemple}\ifx\relax#1\relax\relax\else\textmd{ : \itshape #1}\fi}}
\newtcolorbox{piege}[1][]{boxbase,colback=attncol!6,colframe=attncol,
  title={\textsc{Piège}\ifx\relax#1\relax\relax\else\textmd{ : \itshape #1}\fi}}
\newtcolorbox{remarque}[1][]{boxbase,colback=black!4,colframe=black!45,
  title={\textsc{Remarque}\ifx\relax#1\relax\relax\else\textmd{ : \itshape #1}\fi}}
% --- boîtes EXERCICE / CORRIGÉ (noms EXACTS, auto-comptées) ---
\newtcolorbox[auto counter]{exo}[1][]{boxbase,colback=primary!4,colframe=primary,
  title={\textsc{Exercice}~\thetcbcounter\ifx\relax#1\relax\relax\else\textmd{ --- \itshape #1}\fi}}
\newtcolorbox[auto counter]{corrige}[1][]{boxbase,colback=excol!4,colframe=excol,
  title={\textsc{Corrigé}~\thetcbcounter\ifx\relax#1\relax\relax\else\textmd{ --- \itshape #1}\fi}}
\newcommand{\R}{\mathbb{R}}\newcommand{\N}{\mathbb{N}}\newcommand{\Z}{\mathbb{Z}}\newcommand{\ds}{\displaystyle}
% (un \usepackage{fancyhdr}+en-tête est possible ; il est ignoré côté web)
% =========================================================================
\begin{document}
% THEME_TITLE: Les proportions stœchiométriques et le tableau d'avancement
\sisetup{per-mode=symbol}

\begin{center}
{\Large\bfseries\color{primarydark} Thème 3 — Proportions stœchiométriques et avancement}
\end{center}

\medskip
Une fois l'équation pondérée et les données en moles, les coefficients relient toutes les espèces
entre elles : ils fixent les \textbf{proportions}. Calculer un bilan, c'est appliquer une
\emph{règle de trois} avec ces coefficients — rien de plus.

\begin{cle}[proportionnalité aux coefficients]
Pour $a\,\ce{A}+b\,\ce{B}\ \rightarrow\ c\,\ce{C}+d\,\ce{D}$, les quantités \emph{consommées} et
\emph{formées} sont proportionnelles aux coefficients :
\[
\boxed{\;\dfrac{|\Delta n_{\ce{A}}|}{a}=\dfrac{|\Delta n_{\ce{B}}|}{b}
      =\dfrac{\Delta n_{\ce{C}}}{c}=\dfrac{\Delta n_{\ce{D}}}{d}=\xi\;}
\]
Ce rapport commun est l'\textbf{avancement} $\xi$ (en \si{\mole}) : il mesure « combien de fois » la
réaction s'est produite.
\end{cle}

\begin{methode}[la règle de trois stœchiométrique]
Pour trouver la quantité d'une espèce \ce{Y} à partir d'une espèce \ce{X} connue :
\[
n_{\ce{Y}}=n_{\ce{X}}\times\dfrac{\nu_{\ce{Y}}}{\nu_{\ce{X}}}
\qquad(\nu=\text{coefficient stœchiométrique}).
\]
On multiplie par le coefficient de l'espèce \emph{cherchée}, on divise par celui de l'espèce
\emph{connue}. On termine en \textbf{reconvertissant} dans l'unité demandée ($m=nM$, $V=nV_{\mathrm{m}}$).
\end{methode}

\begin{exemplebox}[combien de réactifs pour une quantité visée ?]
Pour $\ce{2 A + 3 B -> C + 6 D}$ (rendement \qty{100}{\percent}), on veut \qty{0.10}{\mole} de \ce{C}.
On multiplie par chaque coefficient (celui de \ce{C} vaut $1$) :
\[
n_{\ce{A}}=0,10\times\tfrac{2}{1}=\qty{0.20}{\mole};\quad
n_{\ce{B}}=0,10\times\tfrac{3}{1}=\qty{0.30}{\mole};\quad
n_{\ce{D}}=0,10\times\tfrac{6}{1}=\qty{0.60}{\mole}.
\]
\end{exemplebox}

\begin{cle}[le tableau d'avancement]
Chaque quantité s'écrit à partir de l'état initial et de l'avancement $\xi$ :
\[
\renewcommand{\arraystretch}{1.25}
\begin{array}{l|cccc}
 & a\,\ce{A} & b\,\ce{B} & c\,\ce{C} & d\,\ce{D}\\\hline
\text{initial} & n_{\ce{A},0} & n_{\ce{B},0} & 0 & 0\\
\text{avancement }\xi & n_{\ce{A},0}-a\xi & n_{\ce{B},0}-b\xi & c\xi & d\xi
\end{array}
\]
Les réactifs \textbf{diminuent} ($-\nu\xi$), les produits \textbf{augmentent} ($+\nu\xi$). La réaction
s'arrête quand un réactif s'annule : l'avancement vaut alors $\xi_{\max}$.
\end{cle}

\begin{piege}[deux réflexes qui font perdre des points]
\textbf{(1)} On applique la règle de trois \emph{sur des moles}, jamais sur des grammes ou des litres
bruts : convertir d'abord. \textbf{(2)} Le résultat sort en moles : pour une masse, multiplier par
$M$ ; pour un volume de gaz, par $V_{\mathrm{m}}$.
\end{piege}

\end{document}
